\documentclass[11pt,a4paper]{article}

% ── Codificación, fuentes y lenguaje ───────────────────────────────────────────
\usepackage[utf8]{inputenc}
\usepackage[T1]{fontenc}
\usepackage[default,scale=0.95]{sourcesanspro}
\renewcommand{\familydefault}{\sfdefault}
\usepackage[spanish,es-noshorthands,es-tabla]{babel}

% ── Geometría y espaciado ─────────────────────────────────────────────────────
\usepackage[top=2.2cm, bottom=2.5cm, left=2.2cm, right=2.2cm, headheight=15pt]{geometry}
\usepackage{setspace}
\setstretch{1.18}
\usepackage{parskip}

% ── Matemáticas y unidades ────────────────────────────────────────────────────
\usepackage{amsmath}
\usepackage{amssymb}
\usepackage{siunitx}

% ── Circuitos ─────────────────────────────────────────────────────────────────
\usepackage[european, straightvoltages]{circuitikz}

% ── Tablas ────────────────────────────────────────────────────────────────────
\usepackage{booktabs}
\usepackage{tabularx}
\usepackage{array}
\usepackage{longtable}
\usepackage{colortbl}
\usepackage{enumitem}

% ── Colores, cajas e iconos ───────────────────────────────────────────────────
\usepackage[dvipsnames]{xcolor}
\usepackage{tcolorbox}
\tcbuselibrary{skins, breakable}
\usepackage{fontawesome5}
\usepackage{graphicx}

% ── Listados de código (dark-mode infográfico) ────────────────────────────────
%   Estilo base: fondo oscuro con números de línea. Los bloques lstlisting se
%   envuelven automáticamente en una caja tcolorbox redondeada.
\usepackage{listings}

% ── Secciones con barra de color (infografía) ────────────────────────────────
\usepackage{titlesec}
\titleformat{\section}
  {\normalfont\LARGE\bfseries\color{azulNoche}}
  {\colorbox{cyanNeon}{\color{fondoOscuro}\thesection}}{0.7em}{}
  [\vspace{2pt}\color{cyanNeon}\rule{\linewidth}{1.8pt}]
\titlespacing*{\section}{0pt}{24pt}{10pt}
\titleformat{\subsection}
  {\normalfont\large\bfseries\color{azulNoche}}
  {\textcolor{cyanNeon}{\thesubsection}}{0.7em}{}
  [\vspace{1pt}\color{grisLinea}\rule{\linewidth}{0.6pt}]
\titlespacing*{\subsection}{0pt}{14pt}{6pt}
\titleformat{\subsubsection}
  {\normalfont\normalsize\bfseries\color{azulNoche}}
  {\textcolor{cyanNeon}{\thesubsubsection}}{0.7em}{}
\titlespacing*{\subsubsection}{0pt}{10pt}{4pt}

% ── Cabeceras y pies de página ────────────────────────────────────────────────
\usepackage{fancyhdr}
\usepackage{lastpage}
\pagestyle{fancy}
\fancyhf{}
\renewcommand{\headrulewidth}{0pt}
\renewcommand{\footrulewidth}{0pt}
\fancyfoot[C]{%
  \begin{minipage}{\linewidth}
    \vspace{2pt}
    \color{cyanNeon}\rule{\linewidth}{1.2pt}\\[2pt]
    \centering\color{grisTexto}\footnotesize
    Página \thepage\ de \pageref{LastPage}
  \end{minipage}%
}

% ── Hiperenlaces ──────────────────────────────────────────────────────────────
\usepackage[colorlinks=true, linkcolor=azulNoche, urlcolor=cyanNeon]{hyperref}
\sloppy
\emergencystretch 3em

% ── Paleta de colores — tecnológico dark-mode ──────────────────────────────────
\definecolor{fondoOscuro}{HTML}{0D1B2A}%             Dark navy — fondo banda título
\definecolor{verdeTurquesa}{HTML}{004D40}%           Verde turquesa — bandas de marca
\definecolor{azulNoche}{HTML}{1B3A6B}%               Midnight blue — secciones, marcos
\definecolor{azulMedio}{HTML}{2A5298}%               Azul medio — variante de acento
\definecolor{cyanNeon}{HTML}{00C8E0}%                Cyan eléctrico — acentos primarios
\definecolor{verdeSignal}{HTML}{00B887}%             Verde señal — fortalezas, éxito
\definecolor{naranjaVivo}{HTML}{FF7043}%             Naranja vivo — mejoras, advertencias
\definecolor{rojoAlerta}{HTML}{E53935}%              Rojo alerta — errores
\definecolor{grisPapel}{HTML}{F4F6FA}%               Fondo tarjetas claras
\definecolor{grisLinea}{HTML}{C8D0E0}%               Bordes sutiles
\definecolor{grisTexto}{HTML}{6B7A99}%               Texto secundario
\definecolor{amarilloNota}{HTML}{FFB300}%            Notas / advertencias doradas
\definecolor{codigoFondo}{HTML}{1E2D3D}%             Fondo del bloque de código
\definecolor{codigoComentario}{HTML}{7A8FA6}%        Comentarios
\definecolor{codigoCadena}{HTML}{FF9D5C}%            Cadenas / strings
\definecolor{codigoKeyword}{HTML}{00C8E0}%           Palabras clave
\definecolor{codigoNumero}{HTML}{00E5A0}%            Números

% ── Banda de título: portada infográfica ──────────────────────────────────────
%   Diseño en 2 capas: banda de color (por defecto fondo oscuro) + franja de
%   acento cyan abajo. Uso: \bandaTitulo[color]{título}{subtítulo}.
%   El icono se puede cambiar con \renewcommand{\iconoBanda}{...} (FontAwesome).
\newcommand{\iconoBanda}{microchip}
\newcommand{\bandaTitulo}[3][fondoOscuro]{%
  \begin{tcolorbox}[
    enhanced,
    colback=#1, colframe=#1,
    boxrule=0pt, arc=8pt, outer arc=8pt,
    width=\linewidth,
    top=18pt, bottom=6pt, left=14pt, right=14pt,
    borderline south={3.5pt}{0pt}{cyanNeon},
  ]
    \begin{minipage}[c]{0.07\linewidth}
      \centering
      \textcolor{cyanNeon}{\fontsize{30}{30}\selectfont\faIcon{\iconoBanda}}
    \end{minipage}%
    \hspace{8pt}%
    \begin{minipage}[c]{0.88\linewidth}
      {\fontsize{20}{24}\selectfont\bfseries\color{white}#2}\par\vspace{5pt}%
      \textcolor{cyanNeon!80!white}{\small\faIcon{angle-right}~#3}%
    \end{minipage}
    \vspace{4pt}
  \end{tcolorbox}%
  \vspace{8pt}%
}

% ── Tarjetas de datos (infografía) ────────────────────────────────────────────
\tcbset{
  tarjetaDato/.style={
    enhanced, breakable,
    arc=6pt, outer arc=6pt,
    colback=grisPapel, colframe=azulNoche,
    boxrule=1pt,
    fonttitle=\bfseries\small, coltitle=white,
    attach boxed title to top left={yshift=-3mm, xshift=6mm},
    boxed title style={
      colback=azulNoche, colframe=azulNoche,
      arc=4pt, boxrule=0pt,
    },
    drop fuzzy shadow=azulNoche!30!white,
    top=8pt, bottom=8pt, left=10pt, right=10pt,
  },
  tarjetaIcono/.style={
    enhanced, breakable,
    arc=6pt, outer arc=6pt,
    colback=white, colframe=grisLinea, boxrule=0.8pt,
    fonttitle=\bfseries\small, coltitle=azulNoche,
    borderline west={3pt}{0pt}{cyanNeon},
    drop fuzzy shadow=azulNoche!25!white,
    top=6pt, bottom=6pt, left=10pt, right=10pt,
  },
  cajaTitulo/.style={
    enhanced, breakable,
    arc=6pt, outer arc=6pt,
    colback=azulNoche!10!grisPapel, colframe=azulNoche,
    boxrule=1pt,
    fonttitle=\bfseries, coltitle=white,
    attach boxed title to top left={yshift=-3mm, xshift=6mm},
    boxed title style={
      colback=azulNoche, colframe=azulNoche,
      arc=4pt, boxrule=0pt,
    },
    drop fuzzy shadow=azulNoche!25!white,
    top=8pt, bottom=8pt, left=10pt, right=10pt,
  },
  cajaContenido/.style={
    enhanced, breakable,
    arc=5pt, outer arc=5pt,
    colback=grisPapel, colframe=grisLinea,
    boxrule=0.8pt,
    fonttitle=\bfseries\small, coltitle=azulNoche,
    top=6pt, bottom=6pt, left=10pt, right=10pt,
  },
  cajaFortaleza/.style={
    enhanced, breakable,
    arc=6pt, outer arc=6pt,
    colback=verdeSignal!8!white, colframe=verdeSignal,
    boxrule=1pt,
    borderline west={4pt}{0pt}{verdeSignal},
    fonttitle=\bfseries\small, coltitle=verdeSignal!70!black,
    drop fuzzy shadow=verdeSignal!20!white,
    top=6pt, bottom=6pt, left=10pt, right=10pt,
  },
  cajaMejora/.style={
    enhanced, breakable,
    arc=6pt, outer arc=6pt,
    colback=naranjaVivo!8!white, colframe=naranjaVivo,
    boxrule=1pt,
    borderline west={4pt}{0pt}{naranjaVivo},
    fonttitle=\bfseries\small, coltitle=naranjaVivo!80!black,
    drop fuzzy shadow=naranjaVivo!20!white,
    top=6pt, bottom=6pt, left=10pt, right=10pt,
  },
  cajaRecomendacion/.style={
    enhanced, breakable,
    arc=6pt, outer arc=6pt,
    colback=cyanNeon!6!white, colframe=cyanNeon!60!azulNoche,
    boxrule=1pt,
    borderline west={4pt}{0pt}{cyanNeon},
    fonttitle=\bfseries\small, coltitle=azulNoche,
    drop fuzzy shadow=azulNoche!20!white,
    top=6pt, bottom=6pt, left=10pt, right=10pt,
  },
}

% ── Ícono + texto (encabezados de sección y elementos) ────────────────────────
\newcommand{\iconotexto}[2]{\textcolor{cyanNeon}{\faIcon{#1}}~#2}

% ── Listados de código: estilo dark + auto-envoltura ─────────────────────────
\lstdefinestyle{estiloCodigo}{
    backgroundcolor=\color{codigoFondo},
    basicstyle=\ttfamily\small\color{white},
    commentstyle=\color{codigoComentario}\itshape,
    stringstyle=\color{codigoCadena},
    keywordstyle=\color{codigoKeyword}\bfseries,
    numberstyle=\tiny\color{codigoComentario},
    numbers=left,
    numbersep=10pt,
    showstringspaces=false,
    breaklines=true,
    breakatwhitespace=false,
    tabsize=4,
    frame=none,
    captionpos=b,
    aboveskip=6pt,
    belowskip=4pt,
    xleftmargin=14pt,
    framexleftmargin=14pt,
    literate=
      {á}{{\'a}}1 {é}{{\'e}}1 {í}{{\'i}}1 {ó}{{\'o}}1 {ú}{{\'u}}1
      {Á}{{\'A}}1 {É}{{\'E}}1 {Í}{{\'I}}1 {Ó}{{\'O}}1 {Ú}{{\'U}}1
      {ñ}{{\~n}}1 {Ñ}{{\~N}}1 {ü}{{\"u}}1 {Ü}{{\"U}}1
      {¿}{{?`}}1 {¡}{{!`}}1
}
\lstdefinestyle{estiloPython}{
    style=estiloCodigo,
    language=Python,
}
\lstset{style=estiloCodigo}
\BeforeBeginEnvironment{lstlisting}{%
  \begin{tcolorbox}[
    enhanced, breakable,
    arc=6pt, outer arc=6pt,
    colback=codigoFondo, colframe=azulNoche,
    boxrule=1pt,
    drop fuzzy shadow=azulNoche!25!white,
    top=2pt, bottom=2pt, left=0pt, right=0pt,
  ]%
}
\AfterEndEnvironment{lstlisting}{\end{tcolorbox}}

% ── Cajas de contenido temáticas ──────────────────────────────────────────────
\newtcolorbox{cajaRecuerda}{
    enhanced, breakable,
    arc=6pt, outer arc=6pt,
    colback=azulNoche!10!grisPapel,
    colframe=azulNoche, boxrule=1pt,
    borderline west={4pt}{0pt}{cyanNeon},
    fonttitle=\bfseries\small\color{white},
    title={\faIcon{info-circle}~Recuerda},
    attach boxed title to top left={yshift=-3mm, xshift=6mm},
    boxed title style={colback=azulNoche, arc=4pt, boxrule=0pt},
    drop fuzzy shadow=azulNoche!25!white,
    left=10pt, right=10pt, top=8pt, bottom=8pt,
}

\newtcolorbox{cajaConcepto}{
    enhanced, breakable,
    arc=6pt, outer arc=6pt,
    colback=verdeSignal!8!white,
    colframe=verdeSignal, boxrule=1pt,
    borderline west={4pt}{0pt}{verdeSignal},
    fonttitle=\bfseries\small\color{white},
    title={\faIcon{lightbulb}~Concepto clave},
    attach boxed title to top left={yshift=-3mm, xshift=6mm},
    boxed title style={colback=verdeSignal!80!black, arc=4pt, boxrule=0pt},
    drop fuzzy shadow=verdeSignal!20!white,
    left=10pt, right=10pt, top=8pt, bottom=8pt,
}

\newtcolorbox{cajaEjemplo}{
    enhanced, breakable,
    arc=6pt, outer arc=6pt,
    colback=cyanNeon!5!white,
    colframe=cyanNeon!70!azulNoche, boxrule=1pt,
    borderline west={4pt}{0pt}{cyanNeon},
    fonttitle=\bfseries\small\color{fondoOscuro},
    title={\faIcon{code}~Ejemplo},
    attach boxed title to top left={yshift=-3mm, xshift=6mm},
    boxed title style={colback=cyanNeon, arc=4pt, boxrule=0pt},
    drop fuzzy shadow=azulNoche!20!white,
    left=10pt, right=10pt, top=8pt, bottom=8pt,
}

\newtcolorbox{cajaEjercicio}{
    enhanced, breakable,
    arc=6pt, outer arc=6pt,
    colback=azulMedio!7!white,
    colframe=azulMedio, boxrule=1pt,
    borderline west={4pt}{0pt}{azulMedio},
    fonttitle=\bfseries\small\color{white},
    title={\faIcon{pencil-alt}~Ejercicio},
    attach boxed title to top left={yshift=-3mm, xshift=6mm},
    boxed title style={colback=azulMedio, arc=4pt, boxrule=0pt},
    drop fuzzy shadow=azulNoche!20!white,
    left=10pt, right=10pt, top=8pt, bottom=8pt,
}

\newtcolorbox{cajaReto}{
    enhanced, breakable,
    arc=6pt, outer arc=6pt,
    colback=naranjaVivo!6!white,
    colframe=naranjaVivo, boxrule=1pt,
    borderline west={4pt}{0pt}{naranjaVivo},
    fonttitle=\bfseries\small\color{white},
    title={\faIcon{fire}~Reto},
    attach boxed title to top left={yshift=-3mm, xshift=6mm},
    boxed title style={colback=naranjaVivo, arc=4pt, boxrule=0pt},
    drop fuzzy shadow=naranjaVivo!20!white,
    left=10pt, right=10pt, top=8pt, bottom=8pt,
}

% ── Indicadores de dificultad ─────────────────────────────────────────────────
\newcommand{\facil}{\textcolor{verdeSignal}{\faIcon{circle}\,\textbf{Fácil}}}
\newcommand{\intermedio}{\textcolor{naranjaVivo}{\faIcon{circle}\,\textbf{Intermedio}}}
\newcommand{\dificil}{\textcolor{rojoAlerta}{\faIcon{circle}\,\textbf{Difícil}}}

% ─────────────────────────────────────────────────────────────────────────────
\renewcommand{\iconoBanda}{atom}
\fancyhead[L]{\footnotesize\color{grisTexto}\faIcon{atom}~LABORATORIO\_II\_FISICA}
\fancyhead[R]{\footnotesize\color{grisTexto}Sílabo}
\hypersetup{pdftitle={Laboratorio II de Física},pdfauthor={Prof. César Rodríguez}
  pdfsubject={LABORATORIO\_II\_FISICA}}

\begin{document}
\bandaTitulo[fondoOscuro]{Laboratorio II de Física}{Electricidad y magnetismo aplicados: instrumentación, circuitos y tratamiento de datos experimentales}

\begin{center}\small\color{grisTexto}Prof. César Rodríguez \textbar{} Universidad de Oriente, Núcleo de Sucre, Departamento de Física \textbar{} 2 de noviembre de 2026\end{center}

\section{\iconotexto{id-card}{Identificación}}

\begin{tabularx}{\linewidth}{@{}>{\raggedright\arraybackslash}p{4.6cm}X@{}}
\toprule
\textbf{Código} & LABORATORIO\_II\_FISICA \\
\textbf{Duración} & 10 semanas \\
\textbf{Grupos} & 3 grupos de trabajo \\
\textbf{Modalidad} & Laboratorio presencial en mesas de tres grupos con roles rotativos, más análisis de datos en Python y entrega semanal por repositorio privado. \\
\textbf{Entrega} & GitHub por grupo, revisado al finalizar cada semana \\
\textbf{Dedicación} & 3 horas de laboratorio + 2 horas de análisis e informe \\
\textbf{Versión} & 1 \\
\bottomrule
\end{tabularx}

\section{\iconotexto{bullseye}{Objetivo general}}

Formar en la medición eléctrica rigurosa: montar circuitos con seguridad, usar el instrumental sin exceder sus rangos, estimar la incertidumbre de cada magnitud, contrastar el dato con el modelo teórico y comunicar el resultado en un informe trazable y reproducible.

\section{\iconotexto{tasks}{Resultados de aprendizaje}}

\begin{enumerate}[leftmargin=1.4em,itemsep=2pt,topsep=3pt]
  \item RA1: Seleccionar el instrumento y el rango adecuados, justificando la elección por resolución, clase y efecto de carga.
  \item RA2: Calcular la incertidumbre de una magnitud directa e indirecta y propagarla hasta el resultado final con cifras significativas coherentes.
  \item RA3: Verificar experimentalmente la ley de Ohm, las leyes de Kirchhoff y los teoremas de Thévenin y Norton.
  \item RA4: Determinar resistencias internas de instrumentos y de fuentes reales por pendiente y por sustitución.
  \item RA5: Caracterizar elementos no lineales y extraer sus parámetros por linealización de un modelo exponencial.
  \item RA6: Medir reactancias, desfases, resonancia, factor de calidad y ancho de banda en régimen senoidal.
  \item RA7: Determinar la constante de tiempo de un circuito RC y contrastarla con el valor nominal.
  \item RA8: Procesar datos en Python, ajustar por mínimos cuadrados, graficar con unidades y documentar el script.
  \item RA9: Redactar un informe con datos crudos inalterados, análisis reproducible y conclusiones contrastadas con el modelo.
\end{enumerate}

\section{\iconotexto{sitemap}{Estructura por fases}}

\begin{tcolorbox}[cajaContenido,title={Fase 1 --- Instrumentación y leyes básicas de los circuitos}]Semanas 1 a 4. Medir bien antes de medir cosas difíciles: instrumentos, incertidumbres, efecto de carga y verificación de las leyes fundamentales.
\end{tcolorbox}

\begin{tcolorbox}[cajaContenido,title={Fase 2 --- Elementos no lineales y régimen senoidal}]Semanas 5 a 8. Salir del resistor óhmico: diodos, reactancias, resonancia y transitorios, con linealización de modelos y medición de fase.
\end{tcolorbox}

\begin{tcolorbox}[cajaContenido,title={Fase 3 --- Integración y cierre}]Semanas 9 a 10. Caracterizar un circuito completo por dos vías independientes y defender el portafolio del semestre.
\end{tcolorbox}

\section{\iconotexto{clipboard-check}{Sistema de evaluación}}

\begin{tabularx}{\linewidth}{@{}X r@{}}
\toprule
Concepto & Peso (\%)\\
\midrule
Prácticas de laboratorio (semanas 1 a 9) & 60,0 \\
Controles de fase (semanas 4 y 8) & 20,0 \\
Evaluación final práctica y portafolio (semana 10) & 20,0 \\
\midrule
\textbf{Total} & \textbf{100,0}\\
\bottomrule
\end{tabularx}

\begin{cajaRecuerda}La escala de calificación es de 10 puntos y la nota mínima aprobatoria es 5. La nota del curso se calcula ponderando los conceptos de la tabla; no se promedian notas de distinta escala.
\end{cajaRecuerda}

\section{\iconotexto{code-branch}{Entrega y revisión semanal}}

\begin{itemize}[leftmargin=1.4em,itemsep=2pt,topsep=3pt]
  \item Plataforma: GitHub.
  \item Modalidad: Un repositorio privado por grupo, con el docente como colaborador de lectura. El repositorio es la entrega; no se reciben informes por correo ni impresos, salvo indicación expresa del docente.
  \item Cadencia: El docente revisa el repositorio al finalizar cada semana de práctica y devuelve observaciones antes de la semana siguiente.
  \item Corte: Domingo a las 23:59 de la semana correspondiente. Lo entregado después del corte se corrige con la penalización prevista en las políticas.
\end{itemize}

\textbf{Estructura del repositorio}

\begin{itemize}[leftmargin=1.4em,itemsep=2pt,topsep=3pt]
  \item README.md con integrantes, roles de la fase y tabla de semanas entregadas.
  \item semana\_NN/informe.pdf con el informe de la práctica.
  \item semana\_NN/datos/ con los CSV crudos, uno por ensayo, sin editar.
  \item semana\_NN/analisis/ con los scripts de Python que reproducen cada figura.
  \item semana\_NN/figuras/ con las gráficas en PDF o PNG generadas por los scripts.
  \item semana\_NN/README.md con lo que se hizo, lo que no se pudo verificar y por qué.
\end{itemize}

\textbf{Reglas de entrega}

\begin{itemize}[leftmargin=1.4em,itemsep=2pt,topsep=3pt]
  \item Los datos crudos no se editan nunca; toda corrección vive en el script de análisis.
  \item Cada cifra del informe debe salir de una tabla de datos o de un script entregado.
  \item Los scripts se ejecutan de punta a punta sin pasos manuales intermedios.
  \item Nada de credenciales ni tokens en el repositorio; las claves van en .env ignorado.
  \item El informe es grupal, pero las conclusiones se firman individualmente.
  \item Un resultado sin incertidumbre no se acepta como resultado.
\end{itemize}

\begin{cajaRecuerda}La entrega se considera válida solo cuando está en el repositorio del grupo antes del corte. Un enlace, una captura suelta o un archivo enviado por mensajería no sustituyen el repositorio.\end{cajaRecuerda}

\section{\iconotexto{server}{Laboratorio}}

\begin{tabularx}{\linewidth}{@{}l l r X@{}}
\toprule
Contenedor & Imagen & Puerto & Propósito\\
\midrule
analisis & \texttt{jupyter/scipy-notebook:latest} & 8888 & Entorno Jupyter con NumPy, SciPy y Matplotlib para ajustar datos, propagar incertidumbres y generar las figuras del informe. \\
\bottomrule
\end{tabularx}

Requiere alrededor de 2 GB de RAM libre y 3 GB de disco. La imagen incluye NumPy, SciPy, Matplotlib y Pandas, suficientes para todos los ajustes del curso.

El laboratorio físico es presencial y obligatorio: los montajes se hacen en la mesa con instrumentos reales. El contenedor es opcional y sirve solo para el análisis de datos cuando el equipo del grupo no tiene Python instalado.

\begin{itemize}[leftmargin=1.4em,itemsep=2pt,topsep=3pt]
  \item El contenedor procesa datos; nunca sustituye una medición real en la mesa.
  \item Se accede en localhost y no se expone a Internet.
  \item Los CSV crudos se montan de solo lectura para impedir edición accidental.
  \item Cada grupo usa su propio equipo; no hay red ni servidores compartidos del curso.
  \item Lo que no se midió no se simula como si se hubiera medido.
\end{itemize}

\section{\iconotexto{balance-scale}{Políticas de uso}}

\begin{itemize}[leftmargin=1.4em,itemsep=2pt,topsep=3pt]
  \item La seguridad eléctrica es previa a cualquier medición; se corta la alimentación antes de modificar un montaje.
  \item Ningún instrumento se conecta sin verificar rango y polaridad, y ningún capacitor electrolítico se conecta invertido.
  \item Los datos crudos son inalterables; toda corrección se documenta en el script de análisis y se explica en el informe.
  \item Todo resultado se expresa en unidades del Sistema Internacional con cifras significativas coherentes con su incertidumbre.
  \item Un resultado sin incertidumbre no se acepta, y una conclusión sin contraste con el modelo tampoco.
  \item El informe es grupal, pero cada integrante firma sus propias conclusiones y debe poder defenderlas en la oral.
  \item La entrega fuera de corte pierde el 20 por ciento de la nota de la práctica por cada día de retraso, hasta tres días.
  \item Se considera falta grave fabricar, alterar o copiar datos, o entregar el informe de otro grupo; se aplica el reglamento disciplinario.
  \item El instrumental prestado se devuelve completo y revisado; un daño por mal uso se reporta de inmediato al docente.
  \item La asistencia al laboratorio es obligatoria; dos inasistencias injustificadas a prácticas invalidan la evaluación de esa fase.
  \item Los repositorios son privados y no contienen datos personales de terceros ni credenciales.
\end{itemize}

\section{\iconotexto{book}{Bibliografía}}

\begin{itemize}[leftmargin=1.4em,itemsep=2pt,topsep=3pt]
  \item Serway, R. A. y Jewett, J. W.. Física para ciencias e ingeniería. Volúmen 2, capítulos de campos eléctricos, corriente continua y circuitos de corriente alterna. Referencia principal para el modelo teórico.
  \item Halliday, D., Resnick, R. y Walker, J.. Fundamentos de física. Volúmen 2. Tratamiento de circuitos, inducción y corriente alterna con énfasis en la interpretación física.
  \item Tipler, P. A. y Mosca, G.. Física para la ciencia y la tecnología. Volúmen 2. Útil para transitorios RC y RLC y para el análisis de potencia.
  \item Boylestad, R. L.. Introducción al análisis de circuitos. Thévenin, Norton, superposición, respuesta en frecuencia y resonancia con el formalismo fasorial completo.
  \item Sedra, A. S. y Smith, K. C.. Microelectrónica. Modelo del diodo, ecuación de Shockley y diodos Zener y emisores de luz.
  \item Taylor, J. R.. Introducción al análisis de errores. Base del tratamiento de incertidumbres del curso, de la propagación y de los ajustes por mínimos cuadrados.
  \item JCGM. Guía para la expresión de la incertidumbre de medida (GUM, JCGM 100:2008). Norma de referencia para la clasificación de incertidumbres tipo A y tipo B y para el cálculo de la incertidumbre combinada.
\end{itemize}
\end{document}
